\subsection{Shadows, the general graded Ext calculation, and derived recovery}

We use genuine nonnegative weight-graded modules. Put
\[
 S=\Gamma_k(T)=\bigoplus_{n\geq0}ke_n,
 \qquad e_ae_b=\binom{a+b}{a}e_{a+b},
 \qquad\mathcal G=\operatorname{Gr}_{\geq0}(S).
\]
All morphisms in $\mathcal G$ have weight zero. We write
$S[-i]_n=S_{n-i}$, so its free generator has weight $i$.
The cochain grading is independent of this weight grading.

\begin{theorem}[D001: \status{Proved}; exact coefficient extraction]
\label{sa08:D001}
For every commutative ring $k$, define $F(M)=M\otimes_kS$ and
$\kappa(P)=P_0$. There are adjunctions
\[
 L\dashv F\dashv\kappa,
 \qquad LF=\kappa F=\mathrm{Id},
\]
in which all three functors are exact. Explicitly,
\[
 L(P)=\left(\bigoplus_{n\geq0}P_n\right)\Big/
 \left\langle [e_ax]_{n+a}-\binom{n+a}{n}[x]_n
       :x\in P_n,\ a,n\geq0\right\rangle .
\]
\end{theorem}
\begin{proof}
A map from this quotient to $M$ is a family $f_n:P_n\to M$ with
$f_{n+a}(e_ax)=\binom{n+a}{n}f_n(x)$, exactly the condition defining
a graded $S$-map $P\to F(M)$. Thus $L\dashv F$. The other
adjunction is induction and weight-zero evaluation. The latter
functor and $F$ are exact because each $S_n$ is free of rank one.

To prove the additional exactness of $L$, localize at a rational
prime $p$ and put $r_j=p^j-1$. Every $\binom{r_j}{n}$ with
$0\leq n\leq r_j$ is a $p$-unit: modulo $p$ the coefficients of
$(1+z)^{p^j-1}$ are $(-1)^n$. Form the directed system with terms
$P_{r_j}\otimes\mathbb Z_{(p)}$ and transition maps
\[
 x\longmapsto\frac{e_{r_\ell-r_j}x}{\binom{r_\ell}{r_j}}
 \qquad(\ell\geq j).
\]
The multinomial identity
\[
 \binom{r_\ell}{r_j}\binom{r_j}{n}
 =\binom{r_\ell}{n}\binom{r_\ell-n}{r_j-n}
\]
proves composition and independence of the map from $P_n$ to the
colimit represented by $e_{r_j-n}x/\binom{r_j}{n}$. These maps
satisfy precisely the displayed relations for $L$. Conversely the
maps $P_{r_j}\to L(P)$, $x\mapsto[x]_{r_j}$, respect the transition
maps. They give inverse isomorphisms
\[
 L(P)\otimes\mathbb Z_{(p)}
 \simeq\operatorname*{colim}_j(P\otimes\mathbb Z_{(p)})_{r_j}.
\]
Filtered colimits of modules are exact. Localization commutes with
the presentation of $L$, and exactness of abelian groups is detected
at all rational primes. Hence $L$ is exact for every $k$. Finally the
relation with $n=0$ identifies $[me_n]$ with $[me_0]$, proving
$LF(M)=M$.
\end{proof}

\begin{theorem}[D002: \status{Proved}; all finite-tail graded Ext groups]
\label{sa08:D002}
For all $N,i\geq0$, all commutative $k$, and all $k$-modules $M$,
\begin{align*}
 M&\xrightarrow{\ \sim\ }
 \operatorname{Hom}_{\mathcal G}(\tau_NS[-i],F(M)),\\
 \operatorname{Ext}^q_{\mathcal G}(\tau_NS[-i],F(M))&=0
 \qquad(q>0).
\end{align*}
The first map is multiplication by $me_i$. These are also the Ext
groups in the category of all integer-graded $S$-modules.
\end{theorem}
\begin{proof}
The cokernel of $\tau_NS[-i]\to S[-i]$ has bounded weight.
The colimit formula in D001 kills every bounded-weight module, so
$L(\tau_NS[-i])\simeq L(S[-i])$. The latter is $k$: the map is
induced by $e_{n-i}\mapsto\binom ni$, and at $n=r_j$ its scalar is
a $p$-unit for every sufficiently large $j$. Adjunction proves the
stated Hom formula, including its normalization.

An exact left adjoint implies that the right adjoint preserves
injectives; see Stacks, Lemma~12.29.1, tag~015Z. Thus if
$M\to I^\bullet$ is a $k$-injective resolution, then
$F(M)\to F(I^\bullet)$ is an injective resolution in $\mathcal G$.
Applying the Hom formula gives exactly the complex $I^\bullet$,
whose positive cohomology vanishes. A projective resolution of
$\tau_NS[-i]$ by sums of $S[-d]$, $d\geq0$, is projective also in
the integer-graded category; hence its Ext groups agree. The assertion
is not an ungraded Ext assertion and does not include negative $i$.
\end{proof}

\begin{theorem}[D003: \status{Proved}; finite-tail derived unit]
\label{sa08:D003}
Let $\mathcal G_N$ be the modules supported in weights at least $N$.
The exact truncation $\tau_N:\mathcal G\to\mathcal G_N$ has right
adjoint
\[
 (\rho_NP)_i
 =\operatorname{Hom}_{\mathcal G_N}(\tau_NS[-i],P),\quad i\geq0.
\]
Put $\sigma_N=\tau_NF$ and $\psi_N=\kappa\rho_N$. Then
$\sigma_N\dashv\psi_N$, and the derived unit
\[
 K\longrightarrow\mathbf R\psi_N(\sigma_NK)
\]
is an isomorphism for every $K\in D^+(k)$.
\end{theorem}
\begin{proof}
The exact inclusion $j_N:\mathcal G_N\to\mathcal G$ is left
adjoint to $\tau_N$, so truncation preserves injectives. D001 shows
that $F$ does too. The underived unit is the Hom isomorphism of
D002 at $i=0$. Apply it termwise to a bounded-below injective
resolution: $\sigma_N$ takes that resolution to injectives, and
$\psi_N\sigma_NI^\bullet=I^\bullet$. Derived adjunction therefore
also gives full faithfulness of $\sigma_N$ on $D^+(k)$.
\end{proof}

\begin{proposition}[D004: \status{Refuted}; two negligible classes]
\label{sa08:D004}
Modules whose homogeneous elements individually generate bounded-weight
submodules do not in general form a Serre subcategory. Modules of
uniformly bounded weight form a Serre subcategory, but not a localizing
one. These two readings of the manuscript must be distinguished.
\end{proposition}
\begin{proof}
Over $k=\mathbb F_p$ the digit presentation is
\[
 S=\mathbb F_p[x_0,x_1,\ldots]/(x_j^p:j\geq0),
 \qquad\deg(x_j)=p^j.
\]
Let $I=S_{>0}$. Each element of $I/I^2$ is annihilated by $I$, and
$I/I^2$ is the direct sum of one copy of $k$ in each weight $p^j$.
Thus both $I/I^2$ and $k=S/I$ satisfy the elementwise condition.
But in
\[
 0\longrightarrow I/I^2\longrightarrow S/I^2\longrightarrow k
 \longrightarrow0
\]
the class of $1$ generates the unbounded module $S/I^2$. This proves
failure of extension closure. Uniformly bounded modules are closed
under subobjects, quotients and extensions, but
$\bigoplus_{j\geq0}k[-j]$ is unbounded. This is an infinite algebraic
counterexample, proved directly rather than by a finite test.
\end{proof}

\begin{theorem}[D005: \status{Proved}; a precise infinite-shadow quotient]
\label{sa08:D005}
Let $\mathcal T$ be the localizing Serre closure of bounded-weight
modules in $\mathcal G$, and let $Q:\mathcal G\to\mathcal H=
\mathcal G/\mathcal T$ have fully faithful section functor $R$.
The functor $L$ factors through an exact $\overline L:\mathcal H
\to\operatorname{Mod}(k)$. For
\[
 \sigma=QF,\qquad\psi=\kappa R
\]
one has $\overline L\dashv\sigma\dashv\psi$, and
\[
 RQF(M)=F(M),\qquad R^qR(QF(M))=0\quad(q>0),\qquad
 K\simeq\mathbf R\psi\sigma K\quad(K\in D^+(k)).
\]
Thus the repaired infinite-shadow derived unit is an isomorphism.
\end{theorem}
\begin{proof}
The category $\mathcal G$ is a Grothendieck category of modules over
a small graded ring category, so its localizing Serre quotient and
section exist. Exactness and preservation of colimits show that
$\ker L$ is localizing Serre. D001 shows that it contains every
bounded-weight module, hence $\mathcal T$. The exact quotient
universal property gives $\overline L$.

For $V\in\mathcal T$, derived exact adjunction gives
\[
 \operatorname{Ext}^q_{\mathcal G}(V,F(M))
 =\operatorname{Ext}^q_k(LV,M)=0\quad(q\geq0).
\]
This implies saturation. More explicitly, the unit
$F(M)\to RQF(M)$ has torsion kernel and cokernel. Its kernel is
zero by the Hom vanishing. The extension by its cokernel splits by
$\operatorname{Ext}^1$ vanishing. A section object has no torsion
subobject, so the splitting forces the cokernel to vanish as well.
Consequently
\[
 \operatorname{Hom}_{\mathcal H}(QP,QF(M))
 =\operatorname{Hom}_{\mathcal G}(P,F(M))
 =\operatorname{Hom}_k(LP,M),
\]
proving $\overline L\dashv\sigma$. The other adjunction follows
from $Q\dashv R$ and $F\dashv\kappa$. Since $\overline L$ is exact,
$\sigma$ preserves injectives. For a $k$-injective resolution,
$R\sigma I^\bullet=F(I^\bullet)$. The exactness of $F$ proves the
higher vanishing for $R$; applying $\kappa$ proves the derived unit.
The same argument works for bounded-below complexes.
\end{proof}

\begin{theorem}[D006: \status{Proved}; multiplicative derived recovery]
\label{sa08:D006}
For every simplicial set $X$ the genuinely graded complex
\[
 C_{\mathrm{gr}}(X)=\bigoplus_{w\geq0}
       \operatorname{Hom}_{\mathrm{sSet}}(X,\mathcal A_w^\bullet)
\]
has a natural comparison
\[
 c_X:C_{\mathrm{gr}}(X)\longrightarrow F\mathbf R\Gamma(X;k).
\]
It respects the product in the derived category and becomes an
isomorphism after $Q$. In particular, over $\mathbb Z$ the weighted
cochains with products recover $\mathbf R\Gamma(X;\mathbb Z)$ with
its product in $D(\mathbb Z)$ and the ordinary cohomology ring.
\end{theorem}
\begin{proof}
Work in the category $s\mathcal G$ of simplicial graded modules.
The PD Poincare contraction gives a quasi-isomorphism of DGAs
$F_s(k)\to\mathcal A^\bullet$. Define $\Gamma_X$ weightwise by
$(\Gamma_XV)_w=\operatorname{Hom}_{\mathrm{sSet}}(X,V_w)$. It is
right adjoint to the exact functor $P\mapsto S[X]\otimes_SP$,
where $S[X]$ is the free simplicial $S$-module on $X$. Its ordinary-
to-derived comparison therefore gives
\[
 \Gamma_X\mathcal A\longrightarrow\mathbf R\Gamma_X\mathcal A
 \simeq\mathbf R\Gamma_XF_s(k)
 \simeq F\mathbf R\Gamma_X(k).
\]
The last identification follows using injective resolutions: $F_s$
has the exact left adjoint obtained by applying $L$ simplicially,
so preserves injectives, and $\Gamma_XF_s=F\Gamma_X$ termwise.

In weight $w$, this is the resolution edge map, which dimension
shifting identifies as an isomorphism through degree $w$.
The form complex vanishes in degrees greater than $w$. Hence every
cohomology module of $\operatorname{Cone}(c_X)$ has bounded weight.
Exactness of $Q$ makes $Qc_X$ a quasi-isomorphism. This argument
works weightwise for arbitrary infinite $X$.

We verify products as derived morphisms. The diagonal
$S[X]\to S[X]\otimes_SS[X]$, $[x]\mapsto[x]\otimes[x]$, defines
\[
 \mathbf R\Gamma_X(U)\otimes_S^{\mathbf L}\mathbf R\Gamma_X(V)
 \longrightarrow\mathbf R\Gamma_X(U\otimes_S^{\mathbf L}V).
\]
Its adjoint is the diagonal followed by the two derived evaluation
maps. The comparison from ordinary sections is defined by the same
evaluation; the two multiplicativity composites have the identical
adjoint and thus agree. Coassociativity and the counit give the
associativity and unit compatibilities. Each $\mathcal A^q_n$ is
graded $S$-free. The first-quadrant tensor complex
$\mathcal A\otimes_S\mathcal A$ resolves $S$, by its two homology
spectral sequences. Its DGA multiplication therefore represents the
derived multiplication transported from $S$. This proves the
multiplicative assertion, including the derived product itself.

The quotient tensor and full faithfulness needed to recover this
product are proved, with unbounded enhancements, in
Theorem~\ref{sa08:D006E}. Applying the inverse $\mathbf R\psi$ on
the essential image of $\sigma$ recovers the object and the unique
product. Over $\mathbb Z$, derived tensors of bounded-below
complexes remain bounded below because the global dimension is one.
\end{proof}

\paragraph{Provenance and conventions.}
The prompt's R4 says that the general graded Ext corollary is
unproved. The attached \texttt{MPL\_136\_235-281\_notes.md},
N32--N33 (lines 304--328), already supplies the additional exact
adjoint; N71 (lines 646--654) supplies the corrected localizing
quotient. The attached reading \texttt{136\_282-376.md}, lines
40--50, records these results. We independently proved the
claims, including exactness and derived units; this is a correction
of the reported status, with no priority claim. The literal complex
$\operatorname{Hom}(X,\bigoplus_w\mathcal A_w)$ is different for
infinite $X$, as the main comparison counterexample shows. D006
uses the genuinely weight-graded cochains named in A.2.
